\documentclass[a4paper,10pt]{article}

%%%%%%%%%%%% PREÁMBULO %%%%%%%%%%%%%%%%%%%%%

% Paquetes

\usepackage[utf8]{inputenc}
\usepackage[spanish]{babel}
\usepackage[T1]{fontenc}
\usepackage{array}
\usepackage{amsmath}
\usepackage{geometry}
\geometry{margin=2.5cm}
\usepackage{listings}
\usepackage{graphicx}

% Comandos

% Opciones


%%%%%%%%%%%%%%%%%%%%%%%%%%%%%%%%%%%%%%%%%%%%%%

\begin{document}
	
	%Introducción al documento
	\title{Transcripción de operadores}
	\author{Clara Rodríguez, Alejandro Hernández y Lucía Martínez}
	\date{\today}  
	\maketitle 
	
	\begin{abstract}
		\begin{center}
			Este documento está destinado para la transcripción de operadores C a su equivalente en lenguaje circom.
			
		\end{center}
	\end{abstract}
	
	\section*{Guía de contenidos}
	
	\begin{enumerate}
		\item Operadores Lógicos
		\item Operadores de Comparación
		\item Operadores Aritméticos
	\end{enumerate}

	
	\section {Operadores Lógicos}
		\begin{tabular}{|c|>{\centering\arraybackslash}p{5cm}|p{8cm}|}
			\hline
				\textbf{C} & \textbf{Circom} & \textbf{Comentario} \\
			\hline
				\texttt{!a} & 1 - a & Si a es condición, su negación se consigue con la diferencia entre 1 y ésta \\
			\hline
				\texttt{a \&\& b} & a * b & Si a y b son verdaderos, el resultado de su producto será 1, de lo contrario será 0 \\
			\hline
				\texttt{a || b} & 1 - (1-a)*(1-b) ó a+b - (a*b)& Si a, b o ambos son verdaderos, el resultado de su producto será 1, de lo contrario será 0 \\
			\hline
		\end{tabular}
		
	

	
	\section{Operadores de Comparación}
	
		\begin{tabular}{|c|>{\centering\arraybackslash}p{5cm}|p{8cm}|}
			\hline
				\textbf{C} & \textbf{Circom} & \textbf{Comentario} \\
			\hline
				\texttt{a == b} & \texttt{IsZero()(a - b)} & Si \(a = b\), entonces se cumple (==1) \\
			\hline
				\texttt{a != b} & \texttt{1- IsZero()(a - b)} & Si \(a \neq b\), entonces es falso (==0), por tanto se invierte con la resta \\
			\hline
				\texttt{\(a < b\)} & LessThan(n)(a,b) & Si \(a < b\) es verdadero, entonces se cumple (==1) \\
			\hline
				\texttt{\(a \leq b\)} & LessThanEq(n)(a,b) & Si \(a \leq b\) es verdadero, entonces se cumple (==1) \\
			\hline
				\texttt{\(a > b\)} & 1- LessThanEq(n)(a,b) & Si \(a \leq b\) es falso (== 0), entonces \(a > b\) \\
			\hline
				\texttt{\(a \geq b \)} & 1- LessThan(n)(a,b) & Si \(a < b\) es falso (== 0), entonces \(a \geq b \) \\
			\hline
		\end{tabular}
		
		\smallskip
		\noindent \textit{Nota:} Las funciones \texttt{IsZero()} son circuitos que comparan si un valor es igual a 0.
		
		\noindent \textit{Nota:} Las funciones \texttt{LessThan(n)} y \texttt{LessThanEq(n)} son circuitos que comparan dos valores (a < b) de hasta \( n \) bits.
		
		\section{Operadores Aritméticos}
		
		\noindent Tanto en C como en Circom, los operadores suma (+), resta (-) y multiplicación (*) no difieren. En cambio, en operaciones como la división entera (/) y el módulo (\%) si que encontramos diferencias y por ende es necesaria su transcripción.
		
		\vspace{0.3cm}
		
		\noindent La idea desde la que se parte para encontrar esta equivalencia es de la propiedad fundamental de la división entera, la cual establece:
		
		\[
		\scalebox{1.2}{$a = b \cdot q + r$}
		\]
		
		\vspace{0.3cm}
		
		\subsection*{Operación división}
		
		\subsection*{Operación módulo}
		
		\noindent Partiendo de la fórmula \(a = b * q + r\), podemos llegar al resultado de la operación \(a \% b\) con la variable r, la que representa el resto de una división.
		
		\begin{align}
			a &= b \cdot q + r \\
			r &= a - (b \cdot q) \\
			r &= a - (b \cdot ( a / b) )
		\end{align}
		
		\noindent \textit{Nota: } q representa el resultado de una división, por lo que podemos decir que \( q = a/b \).
		
		\vspace{0.4cm}
		
		\begin{tabular}{|c|>{\centering\arraybackslash}p{5cm}|p{8cm}|}
			\hline
			\textbf{C} & \textbf{Circom} & \textbf{Comentario} \\
			\hline
			\texttt{\(a / b\)} & \texttt{} & \\
			\hline
			\texttt{\(a \% b\)} & \texttt{a - (b * (a/b))} & \(a / b\) hace referencia a la equivalencia que se obtiene en circom para poder realizar la operación \% \\
			\hline
		\end{tabular}
		
	
\end{document}